\phantomsection
\chapter*{İzlenceyle Uyumlu Haftalık Çalışma Rotası}
\markboth{İzlenceyle Uyumlu Haftalık Çalışma Rotası}{İzlenceyle Uyumlu Haftalık Çalışma Rotası}
\addcontentsline{toc}{chapter}{İzlenceyle Uyumlu Haftalık Çalışma Rotası}

Bu rota, IMO 301 İstatistik dersinin haftalık konularını kitabın bölümleriyle eşleştirir.

\begin{table}[H]
\centering
\scriptsize
\renewcommand{\arraystretch}{1.12}
\begin{tabularx}{\textwidth}{@{}>{\centering\arraybackslash}p{0.08\textwidth}>{\raggedright\arraybackslash}X>{\raggedright\arraybackslash}p{0.23\textwidth}@{}}
\toprule
Hafta & Ders konusu & Kitaptaki karşılığı \\
\midrule
1 & İstatistiksel düşünme, araştırma süreci ve olasılıktan istatistiğe geçiş & Bölüm 1 \\
2 & Evren, örneklem, parametre, istatistik, veri türleri, örnekleme ve yanlılık & Bölüm 2 ve 6 \\
3 & Merkez, yayılım, dağılım ve grafiklerin yorumlanması & Bölüm 3 \\
4 & Örnekleme değişkenliği ve örnekleme dağılımları & Bölüm 4 \\
5 & Merkezi Limit Teoremi, standart hata ve MLT'nin çıkarımdaki rolü & Bölüm 5 \\
6 & Nokta tahmini; parametre, tahmin edici ve tahmin & Bölüm 7 \\
7 & Ders dışı veya ara değerlendirme haftası & Bölüm 1--7 tekrarı \\
8 & Güven aralığı tahmini, güven düzeyi ve doğru yorumlama & Bölüm 8 \\
9 & Bootstrap, rastgeleleştirme ve hipotez testi mantığı & Bölüm 9--10 \\
10 & Tip I--Tip II hata, test gücü ve tek örneklem \tstat testi & Bölüm 11 \\
11 & Bağımsız ve eşleştirilmiş örneklem \tstat testleri; uygun test seçimi & Bölüm 12 \\
12 & Kategorik veriler, çapraz tablolar ve ki-kare bağımsızlık testi & Bölüm 13 \\
13 & Korelasyon, basit ve çoklu regresyon, karıştırıcılık ve etkileşim & Bölüm 14--15 \\
14 & Genel sınava hazırlık & Bölüm 16 \\
\bottomrule
\end{tabularx}
\end{table}